\section{Higgs $\leftrightarrow$ Zeta Duality}
\label{sec:dualite_higgs_zeta}
\label{sec:dualite_higgs}

The triple identification of \S\ref{sec:surdetermination_un_huitieme}
leads to an unexpected structural conjecture, formulated as
verrou~K4.

\subsection{The common canonical bifurcator}
\label{sec:bifurcateur}

\begin{conjecture}[K4, Higgs $\leftrightarrow$ zeta duality]
\label{conj:K4}
There exists a canonical PT object $B$, the \emph{bifurcator}, such that:
\begin{itemize}[nosep,leftmargin=*]
\item $m_H / v = \mathrm{scale}(B) = \sper = 1/2$ (breathing mode,
  Higgs mass);
\item $\Delta_N = \mathrm{coupling}(B) = \sper^{2}/2 = 1/8$ (Maslov
  phase, HP-PT sector);
\item $\lambdaH = \mathrm{coupling}(B) = \sper^{2}/2 = 1/8$ (Higgs
  self-coupling).
\end{itemize}
\end{conjecture}

The canonical bifurcator is identified as the spinor projector
$\qplus/\qminus$ of the channel $p = 2$ at the fixed point
$\mustar = 15$. Formally,
\begin{equation}
B \;=\; \Pi_+ : \Hilb \,\to\, \Hilb_{\qplus}, \qquad
\Pi_+ \;=\; \tfrac{1}{2}\big(I - i J_{\mathrm{PT}}\big),
\label{eq:bifurcateur}
\end{equation}
where $J_{\mathrm{PT}}$ is the K\"ahler--Fisher complex structure of
$\mathcal{M}_m^{\mathrm{db}}$ (cf.~\S\ref{sec:K3}) and $\Hilb = \Hilb_{\qplus}
\oplus \Hilb_{\qminus}$ the chiral decomposition. The mirror operator
$\Pi_- = (I + i J_{\mathrm{PT}})/2$ projects onto $\qminus$, and $B$
restricts to the channel $p = 2$ (eigenspace of the binary partition
$\Tsix$).

\subsection{The $p = 2$ channel carries both observables}
\label{sec:canal_p2}

The $p = 2$ channel plays three roles in PT
(\cite{PT_M1}):
\begin{itemize}[nosep,leftmargin=*]
\item info / anti-info operator ($\gammap{2} = 0.867$, the most active);
\item bifurcation $\qplus/\qminus$ (persistence chirality);
\item binary partition $\log_2 m$ (theorem $\Tsix$, GFT).
\end{itemize}

The Higgs lives in the channel $p = 2$
(\cite{PT_M5}): it is neutral under
$\mathrm{SU}(3) \times \mathrm{SU}(2) \times \mathrm{U}(1)$, the
three gauges live in $\{3, 5, 7\}$, so only $p = 2$ is admissible
for $H$.

The Maslov phase also lives in the channel $p = 2$: each corner of
the BK double barrier is a point where the sieve switches
$\qplus \leftrightarrow \qminus$, which is the operational expression
of the binary partition $p = 2$ ($\Tsix$).

\subsection{Explicit derivation by the chain $\Tone$ + $\Ttwo$ + Haar + Weyl}
\label{sec:chaine_t1t2haarweyl}

\begin{theoreme}[Higgs--Maslov identity via the canonical chain]
\label{thm:chaine}
At $\sper = 1/2$,
\begin{equation}
\boxed{\;\dfrac{1}{8} \;=\; \dfrac{\sper^{2}}{2} \;=\; \dfrac{|\lambda_2(T_{30})|}{N_{\mathrm{branches}}\,\qplus/\qminus} \;=\; \dfrac{1/4}{2} \;=\; \lambdaH \;=\; \Delta N_{\mathrm{Maslov}}.\;}
\label{eq:K4}
\end{equation}
\end{theoreme}

\begin{proof}[Sketch, four-step chain]
\begin{description}[nosep,leftmargin=*]
\item[$\Tone$.] $\sper = 1/2$ (theorem, \cite{PT_M1}).
\item[$\Ttwo$.] $|\lambda_2(T_{30})| = \sper^{2} = 1/4$: the second
eigenvalue of the primorial transfer matrix $T_{30}$ equals
$\sper^{2}$ (spectral gap, \cite{PT_M1}).
\item[Haar.] Shift $\Reop(s_\zeta) \to \Reop(s_\zeta) + 1/2$ by the
Jacobian of the multiplicative Haar measure
(\S\ref{sec:haar_jacobien}).
\item[Weyl.] Weyl symmetric ordering divides the state count by $2$
(symmetrisation $(u, p_u) \leftrightarrow (p_u, u)$ of
\S\ref{sec:cellule_planck}), which transforms $\sper^{2}$ into
$\sper^{2}/2$.
\end{description}
The composition of the four steps gives
$|\lambda_2|/2 = \sper^{2}/2 = 1/8$. Identifying the Higgs sector to
$\lambdaH = \sper^{2}/2$ and the HP-PT sector to
$\Delta_N^{\mathrm{Maslov}} = \sper^{2}/2$, one obtains \eqref{eq:K4}.
\end{proof}

The quantity $1/8$ is therefore structurally the $\Ttwo$ spectral
gap divided by the number of $\qplus/\qminus$ branches. The two
readings (Higgs and Maslov) are two names for the same number,
generated by the same mechanism.

\subsection{Isolation tests}
\label{sec:isolation}

No other canonical PT coupling equals $1/8$:

\begin{table}[ht]
\centering
\begin{tabular}{@{}lll@{}}
\toprule
\textbf{PT coupling} & \textbf{Value} & \textbf{$= 1/8$?}\\
\midrule
$\alpha_{\mathrm{EM}}$ & $7.30 \cdot 10^{-3}$ & no\\
$\sin^{2}\theta_W$ & $0.231$ & no\\
$G_F$ (GeV$^{-2}$) & $1.17 \cdot 10^{-5}$ & no\\
$\lambdaH$ & $\mathbf{0.125}$ & \textbf{yes}\\
$\Delta N_{\mathrm{Maslov}}$ & $\mathbf{0.125}$ & \textbf{yes}\\
\bottomrule
\end{tabular}
\caption{Isolation tests: no PT coupling other than $\lambdaH$ and
$\Delta N_{\mathrm{Maslov}}$ equals $1/8$.}
\label{tab:isolation}
\end{table}

Pure numerical coincidence is rejected. The spin hierarchy
$(\sper, \sper^{2}, \sper^{2}/2)$ is dense up to $1/8$ and closes
algebraically beyond ($\sper^{3} = \sper^{2}/2$ at $\sper = 1/2$).

\subsection{Quantitative numerical tests}
\label{sec:tests_higgs}

At tree level, $m_H^{\mathrm{tree}} = v \cdot \sper = 123.11$~GeV
(using $v = 246.22$~GeV). At first PT perturbative order,
\begin{equation}
m_H^{(1)} \;=\; v \cdot \sper \cdot (1 + C_F\,\eps) \;\simeq\; 124.65 \text{ GeV},
\label{eq:mH_1}
\end{equation}
then the full PT-NLO value (with higher-order corrections,
\cite{PT_M5}) reaches
$m_H^{\mathrm{NLO}} = 125.287$~GeV, where the two coefficients are
defined intrinsically in PT:
\begin{itemize}[nosep,leftmargin=*]
\item $C_F = (N_c^{2} - 1)/(2 N_c) = 4/3$ is the Casimir factor of
the fundamental representation of $\mathrm{SU}(N_c)$, with $N_c = 3$
imposed by PT as the only value of $N_c$ for which the fixed-point
equation $\mustar = \sum p_i$ admits a solution on the active primes
(theorem $\Tzero$, \cite{PT_M3}). Geometrically, $C_F$
counts the allowed transitions of the transfer matrix $T_3$ of the
$p = 3$ channel (7 transitions over 9 slots;
\cite{PT_M5}). Status \textbf{[Thm]} for
the derivation $N_c = 3 \Rightarrow C_F = 4/3$.
\item $\eps = \alpha_s\,\sper / (2\pi)$ is the small PT parameter of
the one-loop perturbative expansion, combining the strong coupling
$\alpha_s(m_Z) \simeq 0.1179$ and the bifurcation parameter
$\sper = 1/2$ with the standard $1/(2\pi)$ normalisation of a gluon
loop (\cite{PT_M5}).
Numerically $\eps \simeq 9.38 \cdot 10^{-3}$ and
$C_F\,\eps \simeq 1.25 \cdot 10^{-2}$. Status \textbf{[Der]}
(standard perturbative formula with $\alpha_s$ and $\sper$ taken at
their canonical PT values, with no free parameter).
\end{itemize}

The LHC experimental value is $125.25$~GeV. The first-order error
$m_H^{(1)}$ is $\sim 0.5\,\%$, and the full PT-NLO error
$m_H^{\mathrm{NLO}}$ is $0.0295\,\%$.

\subsection{Final position}
\label{sec:K4_position}

K4 is established as a \textbf{structurally validated strong
conjecture}. Pure numerical coincidence is rejected. The common
mechanism is identified: $\qplus/\qminus$ bifurcation of the $p = 2$
channel. The derivation chain $\Tone \to \Ttwo \to$~Haar~$\to$~Weyl
uses only \textbf{[Thm]} theorems of the corpus. Promotion to an
unconditional theorem requires the strict cohomological closure K3
(currently \textbf{[Thm partial]}, \S\ref{sec:limites_ouvertures}).
